\documentclass[aspectratio=169,10pt]{beamer}
\makeatletter\def\input@path{{../../estilo/}{../estilo/}{estilo/}}\makeatother
\usepackage{devmobile}
\modulo[Flutter]{05}{Widgets e Layouts}{Árvore de widgets, componentes básicos, Row/Column/Stack, restrições de layout, listas, Material 3, temas e telas responsivas}

\begin{document}
\capa

\begin{frame}{Objetivos do módulo}
  \begin{conceito}[Ao final deste módulo você será capaz de]
    \begin{itemize}
      \item Explicar a interface declarativa, as árvores do Flutter e o caminho de um quadro.
      \item Usar textos, ícones, imagens, botões, cartões e campos de entrada.
      \item Montar telas com \cod{Row}, \cod{Column}, \cod{Stack}, \cod{Expanded} e espaçamentos.
      \item Entender a regra de restrições do layout e resolver os erros mais comuns.
      \item Exibir listas e grades eficientes com \cod{ListView.builder} e \cod{GridView}.
      \item Aplicar Material 3, temas claro e escuro, e adaptar a tela ao tamanho do aparelho.
    \end{itemize}
  \end{conceito}
\end{frame}

\sumario

% =============================================================================
%  Fundamentos compartilhados --- Módulo 05: interfaces declarativas
%  Incluído pelas trilhas Kotlin e Flutter com % =============================================================================
%  Fundamentos compartilhados --- Módulo 05: interfaces declarativas
%  Incluído pelas trilhas Kotlin e Flutter com % =============================================================================
%  Fundamentos compartilhados --- Módulo 05: interfaces declarativas
%  Incluído pelas trilhas Kotlin e Flutter com \input{../fundamentos/f05-ui-declarativa}
% =============================================================================
\section{Fundamentos: interfaces declarativas}

\begin{frame}{Interfaces imperativas e declarativas}
  \begin{columns}[T]
    \begin{column}{0.49\textwidth}
      \begin{definicao}[Interface imperativa]
        O programa cria os componentes uma vez e depois os \textbf{altera} passo a passo (\cod{setText}, \cod{setVisibility}). Basta esquecer um deles para a tela mostrar dados errados.
      \end{definicao}
    \end{column}
    \begin{column}{0.49\textwidth}
      \begin{definicao}[Interface declarativa]
        O programa \textbf{descreve} como a tela deve ser para cada estado. Quando o estado muda, o framework gera a nova descrição e aplica apenas as diferenças.
      \end{definicao}
    \end{column}
  \end{columns}
  \vspace{2pt}
  \centering
  \begin{tikzpicture}[node distance=2mm and 6mm,every node/.style={font=\scriptsize}]
    \node[rotulo,anchor=east,text width=1.6cm,align=right] (li) at (0,0) {\textbf{imperativo}};
    \node[caixa,right=3mm of li,minimum width=1.6cm] (e1) {estado muda};
    \node[caixaambar,right=of e1,minimum width=3.2cm] (c1) {seu código acha e altera\\ cada componente afetado};
    \node[caixa,right=of c1,minimum width=2.2cm] (t1) {tela (mutável)};
    \draw[seta] (e1) -- (c1); \draw[seta] (c1) -- (t1);
    \node[rotulo,anchor=east,text width=1.6cm,align=right] (ld) at (0,-1.05) {\textbf{declarativo}};
    \node[caixa,right=3mm of ld,minimum width=1.6cm] (e2) {estado muda};
    \node[caixapetroleo,right=of e2,minimum width=2.1cm] (f2) {UI = f(estado)};
    \node[caixa,right=of f2,minimum width=2.1cm] (d2) {nova descrição};
    \node[caixapetroleo,right=of d2,minimum width=2.6cm] (a2) {framework compara\\ e aplica diferenças};
    \draw[seta] (e2) -- (f2); \draw[seta] (f2) -- (d2); \draw[seta] (d2) -- (a2);
  \end{tikzpicture}
  \raggedright
  \begin{historia}
    Views do Android e UIKit do iOS (2008) são imperativos. O React (2013) popularizou a UI declarativa, depois adotada pelo Flutter (2018), SwiftUI (2019) e Jetpack Compose (2021).
  \end{historia}
\end{frame}

\begin{frame}{Composição: a tela como uma árvore}
  \begin{columns}[c]
    \begin{column}{0.52\textwidth}
      \centering
      \begin{tikzpicture}[level distance=10mm,
          every node/.style={caixa,font=\tiny,minimum height=5mm,inner sep=2pt},
          edge from parent/.style={draw=azulpetroleo,line width=0.6pt},
          level 1/.style={sibling distance=30mm},
          level 2/.style={sibling distance=21mm},
          level 3/.style={sibling distance=10mm}]
        \node[caixaescura,font=\tiny] (raiz) {TelaEquipamentos}
          child {node {BarraSuperior}}
          child {node {ListaEquipamentos}
            child {node[caixaambar,font=\tiny] {CartaoBancada}
              child {node {Ícone}}
              child {node {Textos}}
              child {node {Barra}}}
            child {node[caixaambar,font=\tiny] {CartaoBancada}}};
        \draw[seta] (-2.6,0.2) -- (-2.6,-3.2);
        \draw[setaambar] (-2.2,-3.2) -- (-2.2,0.2);
        \node[rotulo,anchor=north west,text width=6cm,align=left] at (-2.75,-3.45) {\textcolor{azulpetroleo}{$\downarrow$} dados descem por parâmetros\quad \textcolor{ambar}{$\uparrow$} eventos sobem por \emph{callbacks}};
      \end{tikzpicture}
    \end{column}
    \begin{column}{0.46\textwidth}
      \begin{definicao}[Composição]
        Construir interfaces complexas \textbf{combinando} componentes pequenos e independentes, em vez de herdar de componentes grandes. Cada componente recebe dados e devolve uma descrição de UI.
      \end{definicao}
      \begin{analogia}
        Como peças de montar: o cartão é feito de ícone, textos e indicador; a lista, de vários cartões. Trocar uma peça não obriga a refazer as outras.
      \end{analogia}
      \begin{dica}[Um bom componente]
        Faz uma coisa só, recebe tudo por parâmetro, não sabe de onde vêm os dados e pode ser visto isoladamente em uma \emph{preview}.
      \end{dica}
    \end{column}
  \end{columns}
\end{frame}

\begin{frame}{Layout por restrições}
  \begin{columns}[c]
    \begin{column}{0.47\textwidth}
      \centering
      \begin{tikzpicture}[node distance=11mm]
        \node[caixaescura,minimum width=3.4cm] (t) {Tela \quad 360 $\times$ 800};
        \node[caixa,minimum width=3.4cm,below=of t] (c) {Coluna};
        \node[caixaambar,minimum width=3.4cm,below=of c] (k) {Cartão};
        \draw[seta] ([xshift=-10mm]t.south) -- node[rotulo,left,text width=2.1cm,align=right]{largura 0--360\\ altura 0--800} ([xshift=-10mm]c.north);
        \draw[seta] ([xshift=-10mm]c.south) -- node[rotulo,left,text width=2.1cm,align=right]{largura 0--360\\ altura livre ($\infty$)} ([xshift=-10mm]k.north);
        \draw[setaambar] ([xshift=10mm]k.north) -- node[rotulo,right,text width=1.6cm,align=left]{``tenho\\ 360 $\times$ 96''} ([xshift=10mm]c.south);
        \node[rotulo,below=1.5mm of k,text width=4cm] {a coluna põe o cartão em $y = 0$\\ e o próximo em $y = 96$};
      \end{tikzpicture}
    \end{column}
    \begin{column}{0.51\textwidth}
      \begin{definicao}[Restrições (\emph{constraints})]
        Intervalo de tamanhos que o pai aceita para o filho: largura e altura mínimas e máximas. \textbf{Justa} quando mínimo = máximo; \textbf{frouxa} quando o mínimo é zero.
      \end{definicao}
      \begin{conceito}[O algoritmo em três frases]
        \textbf{Restrições descem}: o pai diz os limites. \textbf{Tamanhos sobem}: o filho escolhe um tamanho dentro deles. \textbf{O pai posiciona}: o filho não sabe onde está.
      \end{conceito}
      \begin{pordentro}[Uma passada só]
        Cada nó é medido uma única vez por quadro, o que mantém o layout rápido mesmo com milhares de nós. No Compose, cada \cod{Modifier} ajusta as restrições antes de repassá-las; no Flutter, \cod{Padding} e \cod{SizedBox} fazem esse papel.
      \end{pordentro}
    \end{column}
  \end{columns}
\end{frame}

\begin{frame}{Do estado aos pixels: o caminho de um quadro}
  \centering
  \begin{tikzpicture}[node distance=5mm,every node/.style={font=\scriptsize}]
    \node[caixaambar,minimum width=2cm,minimum height=11mm] (s) {estado\\ muda};
    \node[caixapetroleo,right=of s,minimum width=2.3cm,minimum height=11mm] (c) {\textbf{composição}\\ \tiny o que mostrar};
    \node[caixapetroleo,right=of c,minimum width=2.3cm,minimum height=11mm] (l) {\textbf{layout}\\ \tiny tamanho e posição};
    \node[caixapetroleo,right=of l,minimum width=2.3cm,minimum height=11mm] (d) {\textbf{desenho}\\ \tiny cores, textos, formas};
    \node[caixaescura,right=of d,minimum width=2.3cm,minimum height=11mm] (r) {\textbf{rasterização}\\ \tiny GPU gera os pixels};
    \draw[seta] (s) -- (c); \draw[seta] (c) -- (l); \draw[seta] (l) -- (d); \draw[seta] (d) -- (r);
    \node[rotulo,below=1mm of c] {Compose: \emph{Composition}\\ Flutter: \emph{build}};
    \node[rotulo,below=1mm of l] {\emph{Layout}\\ \emph{layout}};
    \node[rotulo,below=1mm of d] {\emph{Drawing}\\ \emph{paint} e \emph{composite}};
    \node[rotulo,below=1mm of r] {motor gráfico\\ (Skia, Impeller)};
  \end{tikzpicture}
  \raggedright
  \begin{columns}[T]
    \begin{column}{0.49\textwidth}
      \begin{definicao}[Quadro e orçamento de tempo]
        A tela é redesenhada no ritmo do sinal vertical do display (VSync): a 60\,Hz cada quadro tem \textbf{16,7\,ms}; a 120\,Hz, só \textbf{8,3\,ms}. Se o trabalho não couber, o quadro se perde e a animação engasga (\emph{jank}).
      \end{definicao}
    \end{column}
    \begin{column}{0.49\textwidth}
      \begin{pordentro}[Pular fases economiza tempo]
        Uma cor que anima só precisa de desenho; uma posição que muda, de layout e desenho. Ler o estado na fase mais tardia possível evita repetir as anteriores.
      \end{pordentro}
      \begin{dica}
        O \emph{Performance} do Flutter DevTools e o \emph{Layout Inspector} do Android Studio mostram onde o tempo do quadro foi gasto.
      \end{dica}
    \end{column}
  \end{columns}
\end{frame}

\begin{frame}{Por dentro: descrição, instância e renderização}
  \begin{columns}[T]
    \begin{column}{0.56\textwidth}
      \scriptsize
      \begin{tabularx}{\linewidth}{@{}L l l@{}}
        \toprule
        \cabecalho{Papel} & \cabecalho{Flutter} & \cabecalho{Compose} \\
        \midrule
        Descrição leve, recriada a cada mudança & \cod{Widget} & chamada \cod{@Composable} \\
        Instância duradoura, guarda o estado & \cod{Element} + \cod{State} & composição (\cod{remember}) \\
        Mede, posiciona e desenha & \cod{RenderObject} & \cod{LayoutNode} \\
        Informa leitores de tela & árvore de semântica & árvore de semântica \\
        \bottomrule
      \end{tabularx}
      \normalsize
      \begin{analogia}
        O widget é a planta do arquiteto: barata de redesenhar. O elemento é o apartamento construído, onde mora o estado. O objeto de renderização é a estrutura medida e pintada. Mudar a planta não obriga a demolir o prédio.
      \end{analogia}
    \end{column}
    \begin{column}{0.42\textwidth}
      \begin{conceito}[Por que separar]
        Descrições são criadas aos milhares por segundo, então precisam ser baratas. Medir e desenhar é caro, então os objetos de renderização são \textbf{reaproveitados} enquanto possível.
      \end{conceito}
      \begin{pordentro}[Semântica e acessibilidade]
        Em paralelo, o framework monta uma árvore de semântica (rótulos, papéis, ações) que o TalkBack e o VoiceOver leem. Ícone sem rótulo fica mudo para quem não enxerga (Módulo 04).
      \end{pordentro}
    \end{column}
  \end{columns}
\end{frame}

\begin{frame}{Atualizar sem refazer tudo}
  \begin{columns}[c]
    \begin{column}{0.42\textwidth}
      \centering
      \begin{tikzpicture}[level distance=10mm,
          every node/.style={caixa,font=\tiny,minimum height=5mm,minimum width=9mm,inner sep=2pt},
          edge from parent/.style={draw=azulpetroleo,line width=0.6pt},
          level 1/.style={sibling distance=22mm},
          level 2/.style={sibling distance=11mm}]
        \node {Tela}
          child {node {Barra} child {node {Título}}}
          child {node[caixaambarforte] (m) {Medidor}
            child {node[caixaambar] {Valor}}
            child {node[caixaambar] {Barra}}};
        \node[rotulo,below=11mm of m,text width=3.4cm] {o estado \cod{temperatura} mudou: só o \textbf{Medidor} e seus filhos são reconstruídos};
      \end{tikzpicture}
    \end{column}
    \begin{column}{0.56\textwidth}
      \begin{conceito}[Flutter: \emph{rebuild} e reconciliação]
        \cod{setState} marca o elemento como sujo; no próximo quadro, só ele executa \cod{build}. Filhos de mesmo tipo e chave são reaproveitados e só atualizados.
      \end{conceito}
      \begin{conceito}[Compose: recomposição inteligente]
        O Compose sabe quais funções leram cada estado e só executa essas de novo; as que recebem parâmetros iguais aos anteriores são puladas (\emph{skipping}).
      \end{conceito}
      \begin{atencao}[Consequência importante]
        Funções de interface podem rodar muitas vezes, em qualquer ordem: devem ser \textbf{rápidas e sem efeitos colaterais}. Rede e banco ficam em APIs de efeito (Módulo 07).
      \end{atencao}
    \end{column}
  \end{columns}
\end{frame}

\begin{frame}[fragile]{Identidade e chaves}
  \begin{columns}[T]
    \begin{column}{0.47\textwidth}
      \centering
      \begin{tikzpicture}[baseline=(current bounding box.north),every node/.style={font=\tiny}]
        \tikzset{item/.style={caixa,minimum width=1.5cm,minimum height=5mm,inner sep=1pt}}
        \node[rotulo] at (0.75,0.55) {antes};
        \node[item,caixaambar] (a) at (0.75,0) {A \quad expandido};
        \node[item] at (0.75,-0.6) {B};
        \node[item] at (0.75,-1.2) {C};
        \node[rotulo,text width=2cm] at (2.95,0.55) {remove A, \textbf{sem} chave};
        \node[item,caixaambar,draw=vermelhoerro] at (2.95,0) {B \quad expandido};
        \node[item] at (2.95,-0.6) {C};
        \node[rotulo,text width=2cm] at (5.15,0.55) {remove A, \textbf{com} chave};
        \node[item] at (5.15,0) {B};
        \node[item] at (5.15,-0.6) {C};
        \node[rotulo,text width=2cm,text=vermelhoerro] at (2.95,-1.25) {o estado ficou na posição 1};
        \node[rotulo,text width=2cm,text=verdeok] at (5.15,-1.25) {o estado saiu junto com A};
      \end{tikzpicture}
\begin{codigo}[title=Compose]{Kotlin}
items(itens, key = { it.id }) { eq ->
  ItemEquipamento(eq)
}
\end{codigo}
\begin{codigo}[title=Flutter]{Dart}
ItemEquipamento(
    key: ValueKey(eq.id), eq: eq, onTap: abrir)
\end{codigo}
    \end{column}
    \begin{column}{0.51\textwidth}
      \begin{definicao}[Chave (\emph{key})]
        Identificador estável que diz ao framework qual item é qual entre duas versões da árvore. Sem chave, a correspondência é feita pela \textbf{posição}.
      \end{definicao}
      \begin{conceito}[Quando importa]
        Em listas cujos itens têm estado próprio (expandido, campo digitado, animação) e podem ser inseridos, removidos ou reordenados.
      \end{conceito}
      \begin{dica}
        Use um identificador do domínio, como o \cod{id} do equipamento. Nunca o índice da lista nem um valor aleatório criado a cada reconstrução.
      \end{dica}
    \end{column}
  \end{columns}
\end{frame}

\begin{frame}{Temas, \emph{tokens} e contexto implícito}
  \begin{columns}[c]
    \begin{column}{0.42\textwidth}
      \centering
      \begin{tikzpicture}[node distance=7mm and 3mm,every node/.style={font=\tiny}]
        \node[caixaambarforte,minimum width=2.4cm,font=\tiny\bfseries] (t) {Tema};
        \node[caixa,below=of t,minimum width=1.6cm,font=\tiny] (s) {Tela};
        \node[caixa,below left=of s,minimum width=1.3cm,font=\tiny] (l) {Lista};
        \node[caixa,below right=of s,minimum width=1.3cm,font=\tiny] (b) {Botão};
        \node[caixa,below=of l,minimum width=1.3cm,font=\tiny] (c) {Cartão};
        \draw[seta] (t) -- (s); \draw[seta] (s) -- (l); \draw[seta] (s) -- (b); \draw[seta] (l) -- (c);
        \draw[setaambar,dashed] (t.east) -| (b.north);
        \draw[setaambar,dashed] (t.west) -| ([xshift=-3mm]c.west) -- (c.west);
        \node[rotulo,text width=4cm,below=2mm of c.south east] {tracejado: componentes leem o tema direto do contexto, sem receber parâmetros};
      \end{tikzpicture}
    \end{column}
    \begin{column}{0.56\textwidth}
      \begin{definicao}[\emph{Tokens} de design]
        Decisões visuais nomeadas pelo \textbf{papel}, não pelo valor: \cod{primary}, \cod{surface}, \cod{titleMedium}. O componente pede ``a cor primária''; o tema decide qual é.
      \end{definicao}
      \begin{pordentro}[Contexto implícito]
        Tema, idioma e tamanho da tela ficam disponíveis para toda uma subárvore: \cod{InheritedWidget} no Flutter (\cod{Theme.of(context)}), \cod{CompositionLocal} no Compose. Quem lê é atualizado quando o valor muda.
      \end{pordentro}
      \begin{dica}
        Com \emph{tokens}, tema escuro e cor dinâmica (Android 12+) funcionam sem mexer nos componentes.
      \end{dica}
    \end{column}
  \end{columns}
\end{frame}

\begin{frame}{Pixels, densidade e telas de todos os tamanhos}
  \begin{columns}[T]
    \begin{column}{0.49\textwidth}
      \centering
      \begin{tikzpicture}[baseline=(current bounding box.north)]
        \draw[cinzamedio,step=0.2] (0,0) grid (1.2,1.2);
        \fill[ambar,opacity=0.7] (0.2,0.2) rectangle (1.0,1.0);
        \node[rotulo] at (0.6,-0.25) {160 dpi: 48\,dp = 48\,px};
        \draw[cinzamedio,step=0.0667] (2.6,0) grid (3.8,1.2);
        \fill[ambar,opacity=0.7] (2.8,0.2) rectangle (3.6,1.0);
        \node[rotulo] at (3.2,-0.25) {480 dpi: 48\,dp = 144\,px};
      \end{tikzpicture}
      \raggedright
      \begin{definicao}[Pixel independente de densidade]
        \cod{px = dp} $\times$ \cod{(dpi} $\div$ \cod{160)}: um botão de 48\,dp tem o mesmo tamanho físico em qualquer tela. No Flutter, é o pixel lógico (fator em \cod{devicePixelRatio}).
      \end{definicao}
      \begin{dica}[Texto]
        Fontes também seguem a escala escolhida pelo usuário: \cod{sp} no Android, \cod{TextScaler} no Flutter.
      \end{dica}
    \end{column}
    \begin{column}{0.49\textwidth}
      \begin{conceito}[Responsivo $\times$ adaptativo]
        \textbf{Responsivo}: o mesmo layout se estica ou reorganiza. \textbf{Adaptativo}: um layout por classe de tamanho de janela (barra inferior no celular, trilho lateral no tablet). Pense em \textbf{espaço disponível}, não em modelo de aparelho.
      \end{conceito}
      \begin{atencao}[Android 16 e telas grandes]
        Com \cod{targetSdk} 36 (exigido pela Play desde agosto de 2026), telas com largura mínima $\geq$ 600\,dp ignoram travas de orientação e de redimensionamento; no API 37 não haverá exceção.
      \end{atencao}
    \end{column}
  \end{columns}
\end{frame}

% =============================================================================
\section{Fundamentos: interfaces declarativas}

\begin{frame}{Interfaces imperativas e declarativas}
  \begin{columns}[T]
    \begin{column}{0.49\textwidth}
      \begin{definicao}[Interface imperativa]
        O programa cria os componentes uma vez e depois os \textbf{altera} passo a passo (\cod{setText}, \cod{setVisibility}). Basta esquecer um deles para a tela mostrar dados errados.
      \end{definicao}
    \end{column}
    \begin{column}{0.49\textwidth}
      \begin{definicao}[Interface declarativa]
        O programa \textbf{descreve} como a tela deve ser para cada estado. Quando o estado muda, o framework gera a nova descrição e aplica apenas as diferenças.
      \end{definicao}
    \end{column}
  \end{columns}
  \vspace{2pt}
  \centering
  \begin{tikzpicture}[node distance=2mm and 6mm,every node/.style={font=\scriptsize}]
    \node[rotulo,anchor=east,text width=1.6cm,align=right] (li) at (0,0) {\textbf{imperativo}};
    \node[caixa,right=3mm of li,minimum width=1.6cm] (e1) {estado muda};
    \node[caixaambar,right=of e1,minimum width=3.2cm] (c1) {seu código acha e altera\\ cada componente afetado};
    \node[caixa,right=of c1,minimum width=2.2cm] (t1) {tela (mutável)};
    \draw[seta] (e1) -- (c1); \draw[seta] (c1) -- (t1);
    \node[rotulo,anchor=east,text width=1.6cm,align=right] (ld) at (0,-1.05) {\textbf{declarativo}};
    \node[caixa,right=3mm of ld,minimum width=1.6cm] (e2) {estado muda};
    \node[caixapetroleo,right=of e2,minimum width=2.1cm] (f2) {UI = f(estado)};
    \node[caixa,right=of f2,minimum width=2.1cm] (d2) {nova descrição};
    \node[caixapetroleo,right=of d2,minimum width=2.6cm] (a2) {framework compara\\ e aplica diferenças};
    \draw[seta] (e2) -- (f2); \draw[seta] (f2) -- (d2); \draw[seta] (d2) -- (a2);
  \end{tikzpicture}
  \raggedright
  \begin{historia}
    Views do Android e UIKit do iOS (2008) são imperativos. O React (2013) popularizou a UI declarativa, depois adotada pelo Flutter (2018), SwiftUI (2019) e Jetpack Compose (2021).
  \end{historia}
\end{frame}

\begin{frame}{Composição: a tela como uma árvore}
  \begin{columns}[c]
    \begin{column}{0.52\textwidth}
      \centering
      \begin{tikzpicture}[level distance=10mm,
          every node/.style={caixa,font=\tiny,minimum height=5mm,inner sep=2pt},
          edge from parent/.style={draw=azulpetroleo,line width=0.6pt},
          level 1/.style={sibling distance=30mm},
          level 2/.style={sibling distance=21mm},
          level 3/.style={sibling distance=10mm}]
        \node[caixaescura,font=\tiny] (raiz) {TelaEquipamentos}
          child {node {BarraSuperior}}
          child {node {ListaEquipamentos}
            child {node[caixaambar,font=\tiny] {CartaoBancada}
              child {node {Ícone}}
              child {node {Textos}}
              child {node {Barra}}}
            child {node[caixaambar,font=\tiny] {CartaoBancada}}};
        \draw[seta] (-2.6,0.2) -- (-2.6,-3.2);
        \draw[setaambar] (-2.2,-3.2) -- (-2.2,0.2);
        \node[rotulo,anchor=north west,text width=6cm,align=left] at (-2.75,-3.45) {\textcolor{azulpetroleo}{$\downarrow$} dados descem por parâmetros\quad \textcolor{ambar}{$\uparrow$} eventos sobem por \emph{callbacks}};
      \end{tikzpicture}
    \end{column}
    \begin{column}{0.46\textwidth}
      \begin{definicao}[Composição]
        Construir interfaces complexas \textbf{combinando} componentes pequenos e independentes, em vez de herdar de componentes grandes. Cada componente recebe dados e devolve uma descrição de UI.
      \end{definicao}
      \begin{analogia}
        Como peças de montar: o cartão é feito de ícone, textos e indicador; a lista, de vários cartões. Trocar uma peça não obriga a refazer as outras.
      \end{analogia}
      \begin{dica}[Um bom componente]
        Faz uma coisa só, recebe tudo por parâmetro, não sabe de onde vêm os dados e pode ser visto isoladamente em uma \emph{preview}.
      \end{dica}
    \end{column}
  \end{columns}
\end{frame}

\begin{frame}{Layout por restrições}
  \begin{columns}[c]
    \begin{column}{0.47\textwidth}
      \centering
      \begin{tikzpicture}[node distance=11mm]
        \node[caixaescura,minimum width=3.4cm] (t) {Tela \quad 360 $\times$ 800};
        \node[caixa,minimum width=3.4cm,below=of t] (c) {Coluna};
        \node[caixaambar,minimum width=3.4cm,below=of c] (k) {Cartão};
        \draw[seta] ([xshift=-10mm]t.south) -- node[rotulo,left,text width=2.1cm,align=right]{largura 0--360\\ altura 0--800} ([xshift=-10mm]c.north);
        \draw[seta] ([xshift=-10mm]c.south) -- node[rotulo,left,text width=2.1cm,align=right]{largura 0--360\\ altura livre ($\infty$)} ([xshift=-10mm]k.north);
        \draw[setaambar] ([xshift=10mm]k.north) -- node[rotulo,right,text width=1.6cm,align=left]{``tenho\\ 360 $\times$ 96''} ([xshift=10mm]c.south);
        \node[rotulo,below=1.5mm of k,text width=4cm] {a coluna põe o cartão em $y = 0$\\ e o próximo em $y = 96$};
      \end{tikzpicture}
    \end{column}
    \begin{column}{0.51\textwidth}
      \begin{definicao}[Restrições (\emph{constraints})]
        Intervalo de tamanhos que o pai aceita para o filho: largura e altura mínimas e máximas. \textbf{Justa} quando mínimo = máximo; \textbf{frouxa} quando o mínimo é zero.
      \end{definicao}
      \begin{conceito}[O algoritmo em três frases]
        \textbf{Restrições descem}: o pai diz os limites. \textbf{Tamanhos sobem}: o filho escolhe um tamanho dentro deles. \textbf{O pai posiciona}: o filho não sabe onde está.
      \end{conceito}
      \begin{pordentro}[Uma passada só]
        Cada nó é medido uma única vez por quadro, o que mantém o layout rápido mesmo com milhares de nós. No Compose, cada \cod{Modifier} ajusta as restrições antes de repassá-las; no Flutter, \cod{Padding} e \cod{SizedBox} fazem esse papel.
      \end{pordentro}
    \end{column}
  \end{columns}
\end{frame}

\begin{frame}{Do estado aos pixels: o caminho de um quadro}
  \centering
  \begin{tikzpicture}[node distance=5mm,every node/.style={font=\scriptsize}]
    \node[caixaambar,minimum width=2cm,minimum height=11mm] (s) {estado\\ muda};
    \node[caixapetroleo,right=of s,minimum width=2.3cm,minimum height=11mm] (c) {\textbf{composição}\\ \tiny o que mostrar};
    \node[caixapetroleo,right=of c,minimum width=2.3cm,minimum height=11mm] (l) {\textbf{layout}\\ \tiny tamanho e posição};
    \node[caixapetroleo,right=of l,minimum width=2.3cm,minimum height=11mm] (d) {\textbf{desenho}\\ \tiny cores, textos, formas};
    \node[caixaescura,right=of d,minimum width=2.3cm,minimum height=11mm] (r) {\textbf{rasterização}\\ \tiny GPU gera os pixels};
    \draw[seta] (s) -- (c); \draw[seta] (c) -- (l); \draw[seta] (l) -- (d); \draw[seta] (d) -- (r);
    \node[rotulo,below=1mm of c] {Compose: \emph{Composition}\\ Flutter: \emph{build}};
    \node[rotulo,below=1mm of l] {\emph{Layout}\\ \emph{layout}};
    \node[rotulo,below=1mm of d] {\emph{Drawing}\\ \emph{paint} e \emph{composite}};
    \node[rotulo,below=1mm of r] {motor gráfico\\ (Skia, Impeller)};
  \end{tikzpicture}
  \raggedright
  \begin{columns}[T]
    \begin{column}{0.49\textwidth}
      \begin{definicao}[Quadro e orçamento de tempo]
        A tela é redesenhada no ritmo do sinal vertical do display (VSync): a 60\,Hz cada quadro tem \textbf{16,7\,ms}; a 120\,Hz, só \textbf{8,3\,ms}. Se o trabalho não couber, o quadro se perde e a animação engasga (\emph{jank}).
      \end{definicao}
    \end{column}
    \begin{column}{0.49\textwidth}
      \begin{pordentro}[Pular fases economiza tempo]
        Uma cor que anima só precisa de desenho; uma posição que muda, de layout e desenho. Ler o estado na fase mais tardia possível evita repetir as anteriores.
      \end{pordentro}
      \begin{dica}
        O \emph{Performance} do Flutter DevTools e o \emph{Layout Inspector} do Android Studio mostram onde o tempo do quadro foi gasto.
      \end{dica}
    \end{column}
  \end{columns}
\end{frame}

\begin{frame}{Por dentro: descrição, instância e renderização}
  \begin{columns}[T]
    \begin{column}{0.56\textwidth}
      \scriptsize
      \begin{tabularx}{\linewidth}{@{}L l l@{}}
        \toprule
        \cabecalho{Papel} & \cabecalho{Flutter} & \cabecalho{Compose} \\
        \midrule
        Descrição leve, recriada a cada mudança & \cod{Widget} & chamada \cod{@Composable} \\
        Instância duradoura, guarda o estado & \cod{Element} + \cod{State} & composição (\cod{remember}) \\
        Mede, posiciona e desenha & \cod{RenderObject} & \cod{LayoutNode} \\
        Informa leitores de tela & árvore de semântica & árvore de semântica \\
        \bottomrule
      \end{tabularx}
      \normalsize
      \begin{analogia}
        O widget é a planta do arquiteto: barata de redesenhar. O elemento é o apartamento construído, onde mora o estado. O objeto de renderização é a estrutura medida e pintada. Mudar a planta não obriga a demolir o prédio.
      \end{analogia}
    \end{column}
    \begin{column}{0.42\textwidth}
      \begin{conceito}[Por que separar]
        Descrições são criadas aos milhares por segundo, então precisam ser baratas. Medir e desenhar é caro, então os objetos de renderização são \textbf{reaproveitados} enquanto possível.
      \end{conceito}
      \begin{pordentro}[Semântica e acessibilidade]
        Em paralelo, o framework monta uma árvore de semântica (rótulos, papéis, ações) que o TalkBack e o VoiceOver leem. Ícone sem rótulo fica mudo para quem não enxerga (Módulo 04).
      \end{pordentro}
    \end{column}
  \end{columns}
\end{frame}

\begin{frame}{Atualizar sem refazer tudo}
  \begin{columns}[c]
    \begin{column}{0.42\textwidth}
      \centering
      \begin{tikzpicture}[level distance=10mm,
          every node/.style={caixa,font=\tiny,minimum height=5mm,minimum width=9mm,inner sep=2pt},
          edge from parent/.style={draw=azulpetroleo,line width=0.6pt},
          level 1/.style={sibling distance=22mm},
          level 2/.style={sibling distance=11mm}]
        \node {Tela}
          child {node {Barra} child {node {Título}}}
          child {node[caixaambarforte] (m) {Medidor}
            child {node[caixaambar] {Valor}}
            child {node[caixaambar] {Barra}}};
        \node[rotulo,below=11mm of m,text width=3.4cm] {o estado \cod{temperatura} mudou: só o \textbf{Medidor} e seus filhos são reconstruídos};
      \end{tikzpicture}
    \end{column}
    \begin{column}{0.56\textwidth}
      \begin{conceito}[Flutter: \emph{rebuild} e reconciliação]
        \cod{setState} marca o elemento como sujo; no próximo quadro, só ele executa \cod{build}. Filhos de mesmo tipo e chave são reaproveitados e só atualizados.
      \end{conceito}
      \begin{conceito}[Compose: recomposição inteligente]
        O Compose sabe quais funções leram cada estado e só executa essas de novo; as que recebem parâmetros iguais aos anteriores são puladas (\emph{skipping}).
      \end{conceito}
      \begin{atencao}[Consequência importante]
        Funções de interface podem rodar muitas vezes, em qualquer ordem: devem ser \textbf{rápidas e sem efeitos colaterais}. Rede e banco ficam em APIs de efeito (Módulo 07).
      \end{atencao}
    \end{column}
  \end{columns}
\end{frame}

\begin{frame}[fragile]{Identidade e chaves}
  \begin{columns}[T]
    \begin{column}{0.47\textwidth}
      \centering
      \begin{tikzpicture}[baseline=(current bounding box.north),every node/.style={font=\tiny}]
        \tikzset{item/.style={caixa,minimum width=1.5cm,minimum height=5mm,inner sep=1pt}}
        \node[rotulo] at (0.75,0.55) {antes};
        \node[item,caixaambar] (a) at (0.75,0) {A \quad expandido};
        \node[item] at (0.75,-0.6) {B};
        \node[item] at (0.75,-1.2) {C};
        \node[rotulo,text width=2cm] at (2.95,0.55) {remove A, \textbf{sem} chave};
        \node[item,caixaambar,draw=vermelhoerro] at (2.95,0) {B \quad expandido};
        \node[item] at (2.95,-0.6) {C};
        \node[rotulo,text width=2cm] at (5.15,0.55) {remove A, \textbf{com} chave};
        \node[item] at (5.15,0) {B};
        \node[item] at (5.15,-0.6) {C};
        \node[rotulo,text width=2cm,text=vermelhoerro] at (2.95,-1.25) {o estado ficou na posição 1};
        \node[rotulo,text width=2cm,text=verdeok] at (5.15,-1.25) {o estado saiu junto com A};
      \end{tikzpicture}
\begin{codigo}[title=Compose]{Kotlin}
items(itens, key = { it.id }) { eq ->
  ItemEquipamento(eq)
}
\end{codigo}
\begin{codigo}[title=Flutter]{Dart}
ItemEquipamento(
    key: ValueKey(eq.id), eq: eq, onTap: abrir)
\end{codigo}
    \end{column}
    \begin{column}{0.51\textwidth}
      \begin{definicao}[Chave (\emph{key})]
        Identificador estável que diz ao framework qual item é qual entre duas versões da árvore. Sem chave, a correspondência é feita pela \textbf{posição}.
      \end{definicao}
      \begin{conceito}[Quando importa]
        Em listas cujos itens têm estado próprio (expandido, campo digitado, animação) e podem ser inseridos, removidos ou reordenados.
      \end{conceito}
      \begin{dica}
        Use um identificador do domínio, como o \cod{id} do equipamento. Nunca o índice da lista nem um valor aleatório criado a cada reconstrução.
      \end{dica}
    \end{column}
  \end{columns}
\end{frame}

\begin{frame}{Temas, \emph{tokens} e contexto implícito}
  \begin{columns}[c]
    \begin{column}{0.42\textwidth}
      \centering
      \begin{tikzpicture}[node distance=7mm and 3mm,every node/.style={font=\tiny}]
        \node[caixaambarforte,minimum width=2.4cm,font=\tiny\bfseries] (t) {Tema};
        \node[caixa,below=of t,minimum width=1.6cm,font=\tiny] (s) {Tela};
        \node[caixa,below left=of s,minimum width=1.3cm,font=\tiny] (l) {Lista};
        \node[caixa,below right=of s,minimum width=1.3cm,font=\tiny] (b) {Botão};
        \node[caixa,below=of l,minimum width=1.3cm,font=\tiny] (c) {Cartão};
        \draw[seta] (t) -- (s); \draw[seta] (s) -- (l); \draw[seta] (s) -- (b); \draw[seta] (l) -- (c);
        \draw[setaambar,dashed] (t.east) -| (b.north);
        \draw[setaambar,dashed] (t.west) -| ([xshift=-3mm]c.west) -- (c.west);
        \node[rotulo,text width=4cm,below=2mm of c.south east] {tracejado: componentes leem o tema direto do contexto, sem receber parâmetros};
      \end{tikzpicture}
    \end{column}
    \begin{column}{0.56\textwidth}
      \begin{definicao}[\emph{Tokens} de design]
        Decisões visuais nomeadas pelo \textbf{papel}, não pelo valor: \cod{primary}, \cod{surface}, \cod{titleMedium}. O componente pede ``a cor primária''; o tema decide qual é.
      \end{definicao}
      \begin{pordentro}[Contexto implícito]
        Tema, idioma e tamanho da tela ficam disponíveis para toda uma subárvore: \cod{InheritedWidget} no Flutter (\cod{Theme.of(context)}), \cod{CompositionLocal} no Compose. Quem lê é atualizado quando o valor muda.
      \end{pordentro}
      \begin{dica}
        Com \emph{tokens}, tema escuro e cor dinâmica (Android 12+) funcionam sem mexer nos componentes.
      \end{dica}
    \end{column}
  \end{columns}
\end{frame}

\begin{frame}{Pixels, densidade e telas de todos os tamanhos}
  \begin{columns}[T]
    \begin{column}{0.49\textwidth}
      \centering
      \begin{tikzpicture}[baseline=(current bounding box.north)]
        \draw[cinzamedio,step=0.2] (0,0) grid (1.2,1.2);
        \fill[ambar,opacity=0.7] (0.2,0.2) rectangle (1.0,1.0);
        \node[rotulo] at (0.6,-0.25) {160 dpi: 48\,dp = 48\,px};
        \draw[cinzamedio,step=0.0667] (2.6,0) grid (3.8,1.2);
        \fill[ambar,opacity=0.7] (2.8,0.2) rectangle (3.6,1.0);
        \node[rotulo] at (3.2,-0.25) {480 dpi: 48\,dp = 144\,px};
      \end{tikzpicture}
      \raggedright
      \begin{definicao}[Pixel independente de densidade]
        \cod{px = dp} $\times$ \cod{(dpi} $\div$ \cod{160)}: um botão de 48\,dp tem o mesmo tamanho físico em qualquer tela. No Flutter, é o pixel lógico (fator em \cod{devicePixelRatio}).
      \end{definicao}
      \begin{dica}[Texto]
        Fontes também seguem a escala escolhida pelo usuário: \cod{sp} no Android, \cod{TextScaler} no Flutter.
      \end{dica}
    \end{column}
    \begin{column}{0.49\textwidth}
      \begin{conceito}[Responsivo $\times$ adaptativo]
        \textbf{Responsivo}: o mesmo layout se estica ou reorganiza. \textbf{Adaptativo}: um layout por classe de tamanho de janela (barra inferior no celular, trilho lateral no tablet). Pense em \textbf{espaço disponível}, não em modelo de aparelho.
      \end{conceito}
      \begin{atencao}[Android 16 e telas grandes]
        Com \cod{targetSdk} 36 (exigido pela Play desde agosto de 2026), telas com largura mínima $\geq$ 600\,dp ignoram travas de orientação e de redimensionamento; no API 37 não haverá exceção.
      \end{atencao}
    \end{column}
  \end{columns}
\end{frame}

% =============================================================================
\section{Fundamentos: interfaces declarativas}

\begin{frame}{Interfaces imperativas e declarativas}
  \begin{columns}[T]
    \begin{column}{0.49\textwidth}
      \begin{definicao}[Interface imperativa]
        O programa cria os componentes uma vez e depois os \textbf{altera} passo a passo (\cod{setText}, \cod{setVisibility}). Basta esquecer um deles para a tela mostrar dados errados.
      \end{definicao}
    \end{column}
    \begin{column}{0.49\textwidth}
      \begin{definicao}[Interface declarativa]
        O programa \textbf{descreve} como a tela deve ser para cada estado. Quando o estado muda, o framework gera a nova descrição e aplica apenas as diferenças.
      \end{definicao}
    \end{column}
  \end{columns}
  \vspace{2pt}
  \centering
  \begin{tikzpicture}[node distance=2mm and 6mm,every node/.style={font=\scriptsize}]
    \node[rotulo,anchor=east,text width=1.6cm,align=right] (li) at (0,0) {\textbf{imperativo}};
    \node[caixa,right=3mm of li,minimum width=1.6cm] (e1) {estado muda};
    \node[caixaambar,right=of e1,minimum width=3.2cm] (c1) {seu código acha e altera\\ cada componente afetado};
    \node[caixa,right=of c1,minimum width=2.2cm] (t1) {tela (mutável)};
    \draw[seta] (e1) -- (c1); \draw[seta] (c1) -- (t1);
    \node[rotulo,anchor=east,text width=1.6cm,align=right] (ld) at (0,-1.05) {\textbf{declarativo}};
    \node[caixa,right=3mm of ld,minimum width=1.6cm] (e2) {estado muda};
    \node[caixapetroleo,right=of e2,minimum width=2.1cm] (f2) {UI = f(estado)};
    \node[caixa,right=of f2,minimum width=2.1cm] (d2) {nova descrição};
    \node[caixapetroleo,right=of d2,minimum width=2.6cm] (a2) {framework compara\\ e aplica diferenças};
    \draw[seta] (e2) -- (f2); \draw[seta] (f2) -- (d2); \draw[seta] (d2) -- (a2);
  \end{tikzpicture}
  \raggedright
  \begin{historia}
    Views do Android e UIKit do iOS (2008) são imperativos. O React (2013) popularizou a UI declarativa, depois adotada pelo Flutter (2018), SwiftUI (2019) e Jetpack Compose (2021).
  \end{historia}
\end{frame}

\begin{frame}{Composição: a tela como uma árvore}
  \begin{columns}[c]
    \begin{column}{0.52\textwidth}
      \centering
      \begin{tikzpicture}[level distance=10mm,
          every node/.style={caixa,font=\tiny,minimum height=5mm,inner sep=2pt},
          edge from parent/.style={draw=azulpetroleo,line width=0.6pt},
          level 1/.style={sibling distance=30mm},
          level 2/.style={sibling distance=21mm},
          level 3/.style={sibling distance=10mm}]
        \node[caixaescura,font=\tiny] (raiz) {TelaEquipamentos}
          child {node {BarraSuperior}}
          child {node {ListaEquipamentos}
            child {node[caixaambar,font=\tiny] {CartaoBancada}
              child {node {Ícone}}
              child {node {Textos}}
              child {node {Barra}}}
            child {node[caixaambar,font=\tiny] {CartaoBancada}}};
        \draw[seta] (-2.6,0.2) -- (-2.6,-3.2);
        \draw[setaambar] (-2.2,-3.2) -- (-2.2,0.2);
        \node[rotulo,anchor=north west,text width=6cm,align=left] at (-2.75,-3.45) {\textcolor{azulpetroleo}{$\downarrow$} dados descem por parâmetros\quad \textcolor{ambar}{$\uparrow$} eventos sobem por \emph{callbacks}};
      \end{tikzpicture}
    \end{column}
    \begin{column}{0.46\textwidth}
      \begin{definicao}[Composição]
        Construir interfaces complexas \textbf{combinando} componentes pequenos e independentes, em vez de herdar de componentes grandes. Cada componente recebe dados e devolve uma descrição de UI.
      \end{definicao}
      \begin{analogia}
        Como peças de montar: o cartão é feito de ícone, textos e indicador; a lista, de vários cartões. Trocar uma peça não obriga a refazer as outras.
      \end{analogia}
      \begin{dica}[Um bom componente]
        Faz uma coisa só, recebe tudo por parâmetro, não sabe de onde vêm os dados e pode ser visto isoladamente em uma \emph{preview}.
      \end{dica}
    \end{column}
  \end{columns}
\end{frame}

\begin{frame}{Layout por restrições}
  \begin{columns}[c]
    \begin{column}{0.47\textwidth}
      \centering
      \begin{tikzpicture}[node distance=11mm]
        \node[caixaescura,minimum width=3.4cm] (t) {Tela \quad 360 $\times$ 800};
        \node[caixa,minimum width=3.4cm,below=of t] (c) {Coluna};
        \node[caixaambar,minimum width=3.4cm,below=of c] (k) {Cartão};
        \draw[seta] ([xshift=-10mm]t.south) -- node[rotulo,left,text width=2.1cm,align=right]{largura 0--360\\ altura 0--800} ([xshift=-10mm]c.north);
        \draw[seta] ([xshift=-10mm]c.south) -- node[rotulo,left,text width=2.1cm,align=right]{largura 0--360\\ altura livre ($\infty$)} ([xshift=-10mm]k.north);
        \draw[setaambar] ([xshift=10mm]k.north) -- node[rotulo,right,text width=1.6cm,align=left]{``tenho\\ 360 $\times$ 96''} ([xshift=10mm]c.south);
        \node[rotulo,below=1.5mm of k,text width=4cm] {a coluna põe o cartão em $y = 0$\\ e o próximo em $y = 96$};
      \end{tikzpicture}
    \end{column}
    \begin{column}{0.51\textwidth}
      \begin{definicao}[Restrições (\emph{constraints})]
        Intervalo de tamanhos que o pai aceita para o filho: largura e altura mínimas e máximas. \textbf{Justa} quando mínimo = máximo; \textbf{frouxa} quando o mínimo é zero.
      \end{definicao}
      \begin{conceito}[O algoritmo em três frases]
        \textbf{Restrições descem}: o pai diz os limites. \textbf{Tamanhos sobem}: o filho escolhe um tamanho dentro deles. \textbf{O pai posiciona}: o filho não sabe onde está.
      \end{conceito}
      \begin{pordentro}[Uma passada só]
        Cada nó é medido uma única vez por quadro, o que mantém o layout rápido mesmo com milhares de nós. No Compose, cada \cod{Modifier} ajusta as restrições antes de repassá-las; no Flutter, \cod{Padding} e \cod{SizedBox} fazem esse papel.
      \end{pordentro}
    \end{column}
  \end{columns}
\end{frame}

\begin{frame}{Do estado aos pixels: o caminho de um quadro}
  \centering
  \begin{tikzpicture}[node distance=5mm,every node/.style={font=\scriptsize}]
    \node[caixaambar,minimum width=2cm,minimum height=11mm] (s) {estado\\ muda};
    \node[caixapetroleo,right=of s,minimum width=2.3cm,minimum height=11mm] (c) {\textbf{composição}\\ \tiny o que mostrar};
    \node[caixapetroleo,right=of c,minimum width=2.3cm,minimum height=11mm] (l) {\textbf{layout}\\ \tiny tamanho e posição};
    \node[caixapetroleo,right=of l,minimum width=2.3cm,minimum height=11mm] (d) {\textbf{desenho}\\ \tiny cores, textos, formas};
    \node[caixaescura,right=of d,minimum width=2.3cm,minimum height=11mm] (r) {\textbf{rasterização}\\ \tiny GPU gera os pixels};
    \draw[seta] (s) -- (c); \draw[seta] (c) -- (l); \draw[seta] (l) -- (d); \draw[seta] (d) -- (r);
    \node[rotulo,below=1mm of c] {Compose: \emph{Composition}\\ Flutter: \emph{build}};
    \node[rotulo,below=1mm of l] {\emph{Layout}\\ \emph{layout}};
    \node[rotulo,below=1mm of d] {\emph{Drawing}\\ \emph{paint} e \emph{composite}};
    \node[rotulo,below=1mm of r] {motor gráfico\\ (Skia, Impeller)};
  \end{tikzpicture}
  \raggedright
  \begin{columns}[T]
    \begin{column}{0.49\textwidth}
      \begin{definicao}[Quadro e orçamento de tempo]
        A tela é redesenhada no ritmo do sinal vertical do display (VSync): a 60\,Hz cada quadro tem \textbf{16,7\,ms}; a 120\,Hz, só \textbf{8,3\,ms}. Se o trabalho não couber, o quadro se perde e a animação engasga (\emph{jank}).
      \end{definicao}
    \end{column}
    \begin{column}{0.49\textwidth}
      \begin{pordentro}[Pular fases economiza tempo]
        Uma cor que anima só precisa de desenho; uma posição que muda, de layout e desenho. Ler o estado na fase mais tardia possível evita repetir as anteriores.
      \end{pordentro}
      \begin{dica}
        O \emph{Performance} do Flutter DevTools e o \emph{Layout Inspector} do Android Studio mostram onde o tempo do quadro foi gasto.
      \end{dica}
    \end{column}
  \end{columns}
\end{frame}

\begin{frame}{Por dentro: descrição, instância e renderização}
  \begin{columns}[T]
    \begin{column}{0.56\textwidth}
      \scriptsize
      \begin{tabularx}{\linewidth}{@{}L l l@{}}
        \toprule
        \cabecalho{Papel} & \cabecalho{Flutter} & \cabecalho{Compose} \\
        \midrule
        Descrição leve, recriada a cada mudança & \cod{Widget} & chamada \cod{@Composable} \\
        Instância duradoura, guarda o estado & \cod{Element} + \cod{State} & composição (\cod{remember}) \\
        Mede, posiciona e desenha & \cod{RenderObject} & \cod{LayoutNode} \\
        Informa leitores de tela & árvore de semântica & árvore de semântica \\
        \bottomrule
      \end{tabularx}
      \normalsize
      \begin{analogia}
        O widget é a planta do arquiteto: barata de redesenhar. O elemento é o apartamento construído, onde mora o estado. O objeto de renderização é a estrutura medida e pintada. Mudar a planta não obriga a demolir o prédio.
      \end{analogia}
    \end{column}
    \begin{column}{0.42\textwidth}
      \begin{conceito}[Por que separar]
        Descrições são criadas aos milhares por segundo, então precisam ser baratas. Medir e desenhar é caro, então os objetos de renderização são \textbf{reaproveitados} enquanto possível.
      \end{conceito}
      \begin{pordentro}[Semântica e acessibilidade]
        Em paralelo, o framework monta uma árvore de semântica (rótulos, papéis, ações) que o TalkBack e o VoiceOver leem. Ícone sem rótulo fica mudo para quem não enxerga (Módulo 04).
      \end{pordentro}
    \end{column}
  \end{columns}
\end{frame}

\begin{frame}{Atualizar sem refazer tudo}
  \begin{columns}[c]
    \begin{column}{0.42\textwidth}
      \centering
      \begin{tikzpicture}[level distance=10mm,
          every node/.style={caixa,font=\tiny,minimum height=5mm,minimum width=9mm,inner sep=2pt},
          edge from parent/.style={draw=azulpetroleo,line width=0.6pt},
          level 1/.style={sibling distance=22mm},
          level 2/.style={sibling distance=11mm}]
        \node {Tela}
          child {node {Barra} child {node {Título}}}
          child {node[caixaambarforte] (m) {Medidor}
            child {node[caixaambar] {Valor}}
            child {node[caixaambar] {Barra}}};
        \node[rotulo,below=11mm of m,text width=3.4cm] {o estado \cod{temperatura} mudou: só o \textbf{Medidor} e seus filhos são reconstruídos};
      \end{tikzpicture}
    \end{column}
    \begin{column}{0.56\textwidth}
      \begin{conceito}[Flutter: \emph{rebuild} e reconciliação]
        \cod{setState} marca o elemento como sujo; no próximo quadro, só ele executa \cod{build}. Filhos de mesmo tipo e chave são reaproveitados e só atualizados.
      \end{conceito}
      \begin{conceito}[Compose: recomposição inteligente]
        O Compose sabe quais funções leram cada estado e só executa essas de novo; as que recebem parâmetros iguais aos anteriores são puladas (\emph{skipping}).
      \end{conceito}
      \begin{atencao}[Consequência importante]
        Funções de interface podem rodar muitas vezes, em qualquer ordem: devem ser \textbf{rápidas e sem efeitos colaterais}. Rede e banco ficam em APIs de efeito (Módulo 07).
      \end{atencao}
    \end{column}
  \end{columns}
\end{frame}

\begin{frame}[fragile]{Identidade e chaves}
  \begin{columns}[T]
    \begin{column}{0.47\textwidth}
      \centering
      \begin{tikzpicture}[baseline=(current bounding box.north),every node/.style={font=\tiny}]
        \tikzset{item/.style={caixa,minimum width=1.5cm,minimum height=5mm,inner sep=1pt}}
        \node[rotulo] at (0.75,0.55) {antes};
        \node[item,caixaambar] (a) at (0.75,0) {A \quad expandido};
        \node[item] at (0.75,-0.6) {B};
        \node[item] at (0.75,-1.2) {C};
        \node[rotulo,text width=2cm] at (2.95,0.55) {remove A, \textbf{sem} chave};
        \node[item,caixaambar,draw=vermelhoerro] at (2.95,0) {B \quad expandido};
        \node[item] at (2.95,-0.6) {C};
        \node[rotulo,text width=2cm] at (5.15,0.55) {remove A, \textbf{com} chave};
        \node[item] at (5.15,0) {B};
        \node[item] at (5.15,-0.6) {C};
        \node[rotulo,text width=2cm,text=vermelhoerro] at (2.95,-1.25) {o estado ficou na posição 1};
        \node[rotulo,text width=2cm,text=verdeok] at (5.15,-1.25) {o estado saiu junto com A};
      \end{tikzpicture}
\begin{codigo}[title=Compose]{Kotlin}
items(itens, key = { it.id }) { eq ->
  ItemEquipamento(eq)
}
\end{codigo}
\begin{codigo}[title=Flutter]{Dart}
ItemEquipamento(
    key: ValueKey(eq.id), eq: eq, onTap: abrir)
\end{codigo}
    \end{column}
    \begin{column}{0.51\textwidth}
      \begin{definicao}[Chave (\emph{key})]
        Identificador estável que diz ao framework qual item é qual entre duas versões da árvore. Sem chave, a correspondência é feita pela \textbf{posição}.
      \end{definicao}
      \begin{conceito}[Quando importa]
        Em listas cujos itens têm estado próprio (expandido, campo digitado, animação) e podem ser inseridos, removidos ou reordenados.
      \end{conceito}
      \begin{dica}
        Use um identificador do domínio, como o \cod{id} do equipamento. Nunca o índice da lista nem um valor aleatório criado a cada reconstrução.
      \end{dica}
    \end{column}
  \end{columns}
\end{frame}

\begin{frame}{Temas, \emph{tokens} e contexto implícito}
  \begin{columns}[c]
    \begin{column}{0.42\textwidth}
      \centering
      \begin{tikzpicture}[node distance=7mm and 3mm,every node/.style={font=\tiny}]
        \node[caixaambarforte,minimum width=2.4cm,font=\tiny\bfseries] (t) {Tema};
        \node[caixa,below=of t,minimum width=1.6cm,font=\tiny] (s) {Tela};
        \node[caixa,below left=of s,minimum width=1.3cm,font=\tiny] (l) {Lista};
        \node[caixa,below right=of s,minimum width=1.3cm,font=\tiny] (b) {Botão};
        \node[caixa,below=of l,minimum width=1.3cm,font=\tiny] (c) {Cartão};
        \draw[seta] (t) -- (s); \draw[seta] (s) -- (l); \draw[seta] (s) -- (b); \draw[seta] (l) -- (c);
        \draw[setaambar,dashed] (t.east) -| (b.north);
        \draw[setaambar,dashed] (t.west) -| ([xshift=-3mm]c.west) -- (c.west);
        \node[rotulo,text width=4cm,below=2mm of c.south east] {tracejado: componentes leem o tema direto do contexto, sem receber parâmetros};
      \end{tikzpicture}
    \end{column}
    \begin{column}{0.56\textwidth}
      \begin{definicao}[\emph{Tokens} de design]
        Decisões visuais nomeadas pelo \textbf{papel}, não pelo valor: \cod{primary}, \cod{surface}, \cod{titleMedium}. O componente pede ``a cor primária''; o tema decide qual é.
      \end{definicao}
      \begin{pordentro}[Contexto implícito]
        Tema, idioma e tamanho da tela ficam disponíveis para toda uma subárvore: \cod{InheritedWidget} no Flutter (\cod{Theme.of(context)}), \cod{CompositionLocal} no Compose. Quem lê é atualizado quando o valor muda.
      \end{pordentro}
      \begin{dica}
        Com \emph{tokens}, tema escuro e cor dinâmica (Android 12+) funcionam sem mexer nos componentes.
      \end{dica}
    \end{column}
  \end{columns}
\end{frame}

\begin{frame}{Pixels, densidade e telas de todos os tamanhos}
  \begin{columns}[T]
    \begin{column}{0.49\textwidth}
      \centering
      \begin{tikzpicture}[baseline=(current bounding box.north)]
        \draw[cinzamedio,step=0.2] (0,0) grid (1.2,1.2);
        \fill[ambar,opacity=0.7] (0.2,0.2) rectangle (1.0,1.0);
        \node[rotulo] at (0.6,-0.25) {160 dpi: 48\,dp = 48\,px};
        \draw[cinzamedio,step=0.0667] (2.6,0) grid (3.8,1.2);
        \fill[ambar,opacity=0.7] (2.8,0.2) rectangle (3.6,1.0);
        \node[rotulo] at (3.2,-0.25) {480 dpi: 48\,dp = 144\,px};
      \end{tikzpicture}
      \raggedright
      \begin{definicao}[Pixel independente de densidade]
        \cod{px = dp} $\times$ \cod{(dpi} $\div$ \cod{160)}: um botão de 48\,dp tem o mesmo tamanho físico em qualquer tela. No Flutter, é o pixel lógico (fator em \cod{devicePixelRatio}).
      \end{definicao}
      \begin{dica}[Texto]
        Fontes também seguem a escala escolhida pelo usuário: \cod{sp} no Android, \cod{TextScaler} no Flutter.
      \end{dica}
    \end{column}
    \begin{column}{0.49\textwidth}
      \begin{conceito}[Responsivo $\times$ adaptativo]
        \textbf{Responsivo}: o mesmo layout se estica ou reorganiza. \textbf{Adaptativo}: um layout por classe de tamanho de janela (barra inferior no celular, trilho lateral no tablet). Pense em \textbf{espaço disponível}, não em modelo de aparelho.
      \end{conceito}
      \begin{atencao}[Android 16 e telas grandes]
        Com \cod{targetSdk} 36 (exigido pela Play desde agosto de 2026), telas com largura mínima $\geq$ 600\,dp ignoram travas de orientação e de redimensionamento; no API 37 não haverá exceção.
      \end{atencao}
    \end{column}
  \end{columns}
\end{frame}


% =============================================================================
\section{Tudo é widget}

\begin{frame}{A árvore de widgets}
  \begin{columns}[c]
    \begin{column}{0.5\textwidth}
      \centering
      \begin{tikzpicture}[level distance=9mm,sibling distance=24mm,
          every node/.style={caixa,font=\tiny,minimum height=5mm,inner sep=2pt},
          edge from parent/.style={draw=azulpetroleo,line width=0.6pt}]
        \node[caixaescura] {MaterialApp}
          child {node[caixaescura] {Scaffold}
            child {node {AppBar} child {node {Text}}}
            child {node {Column}
              child {node {Text}}
              child {node[caixaambar] {FilledButton} child {node {Text}}}}};
      \end{tikzpicture}
    \end{column}
    \begin{column}{0.48\textwidth}
      \begin{itemize}
        \item Toda a interface é uma \textbf{árvore de widgets}: layout, estilo, gestos e até espaçamento são widgets.
        \item Um widget é uma \textbf{descrição imutável} de parte da tela.
        \item Quando o estado muda, o Flutter chama \cod{build} de novo e compara a nova descrição com a anterior, redesenhando só o necessário.
      \end{itemize}
      \begin{destaque}
        \textbf{UI = f(estado)}
      \end{destaque}
    \end{column}
  \end{columns}
\end{frame}

\begin{frame}[fragile]{StatelessWidget e StatefulWidget}
  \begin{columns}[T]
    \begin{column}{0.5\textwidth}
\begin{codigo}[title=Sem estado]{Dart}
class CartaoBancada extends StatelessWidget {
  const CartaoBancada({
    super.key,
    required this.numero,
    required this.temperatura,
  });

  final int numero;
  final double temperatura;

  @override
  Widget build(BuildContext context) {
    return Card(
      child: ListTile(
        title: Text('Bancada $numero'),
        subtitle: Text('${temperatura.toStringAsFixed(1)} °C'),
      ),
    );
  }
}
\end{codigo}
    \end{column}
    \begin{column}{0.48\textwidth}
\begin{codigo}[title=Com estado (detalhes no Módulo 07)]{Dart}
class Contador extends StatefulWidget {
  const Contador({super.key});
  @override
  State<Contador> createState() => _ContadorState();
}

class _ContadorState extends State<Contador> {
  int cliques = 0;
  @override
  Widget build(BuildContext context) => FilledButton(
    onPressed: () => setState(() => cliques++),
    child: Text('Cliques: $cliques'),
  );
}
\end{codigo}
      \begin{dica}
        \cod{BuildContext} é a posição do widget na árvore; serve para acessar tema, tamanho da tela e navegação.
      \end{dica}
    \end{column}
  \end{columns}
\end{frame}

% =============================================================================
\section{Widgets básicos}

\begin{frame}[fragile]{Texto, ícones, imagens e botões}
  \begin{columns}[T]
    \begin{column}{0.54\textwidth}
\begin{codigo}[title=Componentes]{Dart}
Text(
  'Bancada 3 --- Osciloscópio',
  style: Theme.of(context).textTheme.titleMedium,
  maxLines: 1,
  overflow: TextOverflow.ellipsis,
),
const Icon(Icons.thermostat, color: Colors.orange),
Image.asset('assets/imagens/osciloscopio.png',
    width: 64, fit: BoxFit.cover),
Image.network(equipamento.fotoUrl, height: 120),

FilledButton(onPressed: reservar,
    child: const Text('Reservar')),
OutlinedButton(onPressed: () {},
    child: const Text('Cancelar')),
IconButton(onPressed: abrirMenu,
    icon: const Icon(Icons.more_vert),
    tooltip: 'Menu'),
\end{codigo}
    \end{column}
    \begin{column}{0.44\textwidth}
      \scriptsize
      \begin{tabular}{@{}l l@{}}
        \toprule
        \cabecalho{Widget} & \cabecalho{Uso} \\
        \midrule
        \cod{FilledButton} & ação principal \\
        \cod{OutlinedButton} & ação secundária \\
        \cod{TextButton} & ação discreta \\
        \cod{IconButton} & ícone clicável \\
        \cod{Card}, \cod{ListTile} & itens de lista \\
        \cod{Chip}, \cod{Badge} & etiquetas e contadores \\
        \cod{CircularProgressIndicator} & carregamento \\
        \bottomrule
      \end{tabular}
      \begin{atencao}[Acessibilidade]
        Dê \cod{tooltip} a botões de ícone e \cod{semanticLabel} a imagens informativas.
      \end{atencao}
      \begin{dica}
        \cod{onPressed: null} desabilita o botão.
      \end{dica}
    \end{column}
  \end{columns}
\end{frame}

\begin{frame}[fragile]{Campos de entrada e seleção}
  \begin{columns}[T]
    \begin{column}{0.54\textwidth}
\begin{codigo}[title=Formulário simples]{Dart}
final _form = GlobalKey<FormState>();
final _descricao = TextEditingController();

Form(
  key: _form,
  child: Column(children: [
    TextFormField(
      controller: _descricao,
      decoration: const InputDecoration(
        labelText: 'Descrição',
        border: OutlineInputBorder()),
      validator: (v) =>
          (v == null || v.isEmpty) ? 'Obrigatório' : null,
    ),
    SwitchListTile(
      title: const Text('Alimentação'),
      value: ligado,
      onChanged: (v) => setState(() => ligado = v)),
  ]),
);
// ao salvar: if (_form.currentState!.validate()) ...
\end{codigo}
    \end{column}
    \begin{column}{0.44\textwidth}
      \begin{conceito}[\cod{TextEditingController}]
        Lê e altera o texto do campo. Crie-o no \cod{State} e libere em \cod{dispose()}.
      \end{conceito}
      \scriptsize
      \begin{tabular}{@{}l l@{}}
        \toprule
        \cabecalho{Widget} & \cabecalho{Valor} \\
        \midrule
        \cod{Switch}, \cod{Checkbox} & booleano \\
        \cod{Radio}, \cod{SegmentedButton} & uma opção \\
        \cod{Slider}, \cod{RangeSlider} & número, faixa \\
        \cod{DropdownMenu} & lista suspensa \\
        \cod{showDatePicker} & data \\
        \bottomrule
      \end{tabular}
      \begin{dica}
        \cod{keyboardType: TextInputType.number} abre o teclado numérico.
      \end{dica}
    \end{column}
  \end{columns}
\end{frame}

% =============================================================================
\section{Layout}

\begin{frame}{Row, Column e Stack}
  \begin{columns}[T]
    \begin{column}{0.32\textwidth}
      \centering
      \begin{tikzpicture}
        \draw[azulpetroleo,thick,rounded corners=2pt] (0,0) rectangle (2.4,2.4);
        \foreach \y/\t in {1.85/A,1.2/B,0.55/C}{\node[caixaescura,minimum width=1.9cm,minimum height=5mm] at (1.2,\y) {\t};}
        \draw[setaambar] (2.65,2.1) -- (2.65,0.3);
      \end{tikzpicture}\\[2pt]
      \textbf{\cod{Column}} --- vertical
    \end{column}
    \begin{column}{0.32\textwidth}
      \centering
      \begin{tikzpicture}
        \draw[azulpetroleo,thick,rounded corners=2pt] (0,0) rectangle (2.4,2.4);
        \foreach \x/\t in {0.42/A,1.2/B,1.98/C}{\node[caixaescura,minimum width=6mm,minimum height=1.9cm] at (\x,1.2) {\t};}
        \draw[setaambar] (0.2,-0.25) -- (2.2,-0.25);
      \end{tikzpicture}\\[2pt]
      \textbf{\cod{Row}} --- horizontal
    \end{column}
    \begin{column}{0.32\textwidth}
      \centering
      \begin{tikzpicture}
        \draw[azulpetroleo,thick,rounded corners=2pt] (0,0) rectangle (2.4,2.4);
        \node[caixa,minimum width=2cm,minimum height=2cm] at (1.2,1.2) {};
        \node[caixaescura,minimum width=1.2cm,minimum height=1cm] at (1.2,1.2) {B};
        \node[caixaambarforte,minimum width=5mm,minimum height=5mm] at (2.0,2.0) {C};
      \end{tikzpicture}\\[2pt]
      \textbf{\cod{Stack}} --- sobreposição
    \end{column}
  \end{columns}
  \vspace{4pt}
  \begin{columns}[T]
    \begin{column}{0.48\textwidth}
      \begin{conceito}[Eixo principal (\cod{mainAxisAlignment})]
        \cod{start}, \cod{center}, \cod{end}, \cod{spaceBetween}, \cod{spaceEvenly}; espaço fixo com \cod{spacing: 8}.
      \end{conceito}
    \end{column}
    \begin{column}{0.48\textwidth}
      \begin{conceito}[Eixo cruzado (\cod{crossAxisAlignment})]
        \cod{start}, \cod{center}, \cod{end}, \cod{stretch}. Em \cod{Stack}, posicione filhos com \cod{Positioned} ou \cod{Align}.
      \end{conceito}
    \end{column}
  \end{columns}
\end{frame}

\begin{frame}[fragile]{Compondo um item de lista}
  \begin{columns}[T]
    \begin{column}{0.58\textwidth}
\begin{codigo}[title=item\_equipamento.dart]{Dart}
class ItemEquipamento extends StatelessWidget {
  const ItemEquipamento(
      {super.key, required this.eq, required this.onTap});
  final Equipamento eq;
  final VoidCallback onTap;

  @override
  Widget build(BuildContext context) => Card(
    child: ListTile(
      onTap: onTap,
      leading: const Icon(Icons.memory),
      title: Text(eq.descricao,
          style: const TextStyle(fontWeight: FontWeight.bold)),
      subtitle: Text(eq.patrimonio),
      trailing: Icon(Icons.circle, size: 12,
          color: eq.disponivel ? Colors.green : Colors.orange),
    ),
  );
}
\end{codigo}
    \end{column}
    \begin{column}{0.4\textwidth}
      \small
      \begin{itemize}
        \item \cod{ListTile}: linha pronta com \cod{leading}, \cod{title}, \cod{subtitle} e \cod{trailing}.
        \item Para layouts próprios: \cod{Row} + \cod{Expanded} (ocupa o espaço restante) + \cod{Padding}.
        \item \cod{InkWell} adiciona toque com efeito de onda a qualquer widget.
      \end{itemize}
      \begin{dica}[Extraia widgets em classes]
        Prefira classes \cod{StatelessWidget} a métodos que devolvem widgets: o Flutter consegue reaproveitar e evitar reconstruções.
      \end{dica}
    \end{column}
  \end{columns}
\end{frame}

\begin{frame}[fragile]{Restrições na prática: erros comuns}
  \begin{columns}[T]
    \begin{column}{0.5\textwidth}
\begin{codigo}[title=Lista dentro de coluna]{Dart}
Column(
  children: [
    const CabecalhoResumo(),
    Expanded(            // ocupa o espaço que sobra
      child: ListView.builder(
        itemCount: itens.length,
        itemBuilder: (_, i) => ItemEquipamento(
            eq: itens[i], onTap: () => abrir(itens[i])),
      ),
    ),
  ],
)
\end{codigo}
      \begin{dica}
        O \textbf{Layout Explorer} do DevTools mostra as restrições de cada widget.
      \end{dica}
    \end{column}
    \begin{column}{0.48\textwidth}
      \begin{erro}[``Vertical viewport was given unbounded height'']
        Uma \cod{ListView} dentro de uma \cod{Column} recebe altura ilimitada e não sabe que altura ter. Envolva-a em \cod{Expanded} ou dê uma altura fixa.
      \end{erro}
      \begin{erro}[``A RenderFlex overflowed by 40 pixels'']
        O conteúdo de uma \cod{Row} não coube (faixa amarela e preta). Use \cod{Expanded}/\cod{Flexible}, \cod{Wrap} ou texto com \cod{overflow}.
      \end{erro}
    \end{column}
  \end{columns}
\end{frame}

% =============================================================================
\section{Listas e grades}

\begin{frame}[fragile]{ListView e GridView}
  \begin{columns}[T]
    \begin{column}{0.54\textwidth}
\begin{codigo}[title=lista\_equipamentos.dart]{Dart}
ListView.separated(
  padding: const EdgeInsets.all(16),
  itemCount: itens.length,
  separatorBuilder: (_, __) => const SizedBox(height: 8),
  itemBuilder: (context, i) => ItemEquipamento(
    key: ValueKey(itens[i].patrimonio),
    eq: itens[i],
    onTap: () => abrir(itens[i]),
  ),
);

GridView.builder(
  gridDelegate:
      const SliverGridDelegateWithMaxCrossAxisExtent(
          maxCrossAxisExtent: 200, childAspectRatio: 1.2),
  itemCount: bancadas.length,
  itemBuilder: (context, i) =>
      CartaoBancada(numero: bancadas[i].id,
          temperatura: bancadas[i].temperatura),
);
\end{codigo}
    \end{column}
    \begin{column}{0.44\textwidth}
      \begin{conceito}[\cod{.builder} = sob demanda]
        Constrói só os itens visíveis. Uma \cod{ListView(children: [\ldots])} com mil itens cria todos de uma vez.
      \end{conceito}
      \begin{dica}[Keys]
        \cod{ValueKey} com um identificador estável preserva o estado de cada item quando a lista é reordenada ou filtrada.
      \end{dica}
      \begin{dica}
        \cod{RefreshIndicator}: ``puxar para atualizar''. \cod{CustomScrollView} + \emph{slivers}: cabeçalhos e listas combinados.
      \end{dica}
    \end{column}
  \end{columns}
\end{frame}

% =============================================================================
\section{Material 3, temas e responsividade}

\begin{frame}[fragile]{Scaffold e tema}
  \begin{columns}[T]
    \begin{column}{0.54\textwidth}
\begin{codigo}[title=app.dart]{Dart}
MaterialApp(
  theme: ThemeData(
    colorScheme: ColorScheme.fromSeed(
        seedColor: const Color(0xFF0D3052)),
  ),
  darkTheme: ThemeData(
    colorScheme: ColorScheme.fromSeed(
        seedColor: const Color(0xFF0D3052),
        brightness: Brightness.dark),
  ),
  themeMode: ThemeMode.system,
  home: Scaffold(
    appBar: AppBar(title: const Text('LabMonitor')),
    floatingActionButton: FloatingActionButton(
        onPressed: novo, child: const Icon(Icons.add)),
    body: const ListaEquipamentos(),
  ),
);
\end{codigo}
    \end{column}
    \begin{column}{0.44\textwidth}
      \begin{conceito}[Use o tema]
        \cod{Theme.of(context).colorScheme.primary} e \cod{textTheme.titleLarge} se adaptam ao tema escuro. Cores fixas como \cod{Colors.red} não.
      \end{conceito}
      \begin{dica}[Cupertino]
        Para visual iOS, há os widgets \cod{Cupertino*} e construtores \cod{.adaptive} (ex.: \cod{Switch.adaptive}), que mudam conforme a plataforma.
      \end{dica}
      \begin{dica}
        O pacote \cod{google\_fonts} aplica fontes do Google Fonts ao \cod{textTheme}.
      \end{dica}
    \end{column}
  \end{columns}
\end{frame}

\begin{frame}[fragile]{Telas responsivas}
  \begin{columns}[T]
    \begin{column}{0.54\textwidth}
\begin{codigo}[title=layout\_adaptativo.dart]{Dart}
LayoutBuilder(builder: (context, c) {
  if (c.maxWidth >= 840) {          // expandido
    return Row(children: [
      const NavigationRail(...),
      const Expanded(child: ListaEDetalhe()),
    ]);
  }
  return const TelaCompacta();      // celular
});

final largura = MediaQuery.sizeOf(context).width;

SafeArea(child: conteudo);  // evita entalhe e barras
\end{codigo}
      \begin{dica}
        Unidades no Flutter são \textbf{pixels lógicos} (como o \cod{dp}): o mesmo número tem o mesmo tamanho físico em telas de densidades diferentes.
      \end{dica}
    \end{column}
    \begin{column}{0.44\textwidth}
      \scriptsize
      \begin{tabularx}{\linewidth}{@{}l l L@{}}
        \toprule
        \cabecalho{Classe} & \cabecalho{Largura} & \cabecalho{Navegação} \\
        \midrule
        compacta & < 600 & \cod{NavigationBar} \\
        média & 600--839 & \cod{NavigationRail} \\
        expandida & $\geq$ 840 & lista e detalhe lado a lado \\
        \bottomrule
      \end{tabularx}
      \begin{atencao}[Alvo de toque]
        Elementos clicáveis com pelo menos \textbf{48$\times$48} pixels lógicos.
      \end{atencao}
      \begin{atencao}[Fontes grandes]
        Teste com a fonte do sistema em 200\,\%: textos não podem ser cortados.
      \end{atencao}
    \end{column}
  \end{columns}
\end{frame}

% =============================================================================
\section{Fechamento}

\begin{frame}{Resumo}
  \begin{resumo}
    \begin{itemize}
      \item A interface é uma árvore de widgets imutáveis; \cod{build} descreve a tela para o estado atual.
      \item \cod{StatelessWidget} só exibe; \cod{StatefulWidget} guarda estado e chama \cod{setState}.
      \item \cod{Row}, \cod{Column}, \cod{Stack}, \cod{Expanded} e \cod{Padding} resolvem quase todo layout.
      \item Restrições descem, tamanhos sobem: listas em colunas precisam de \cod{Expanded}.
      \item \cod{ListView.builder} e \cod{GridView.builder} para listas grandes, com \cod{ValueKey}.
      \item \cod{ColorScheme.fromSeed}, tema escuro e \cod{LayoutBuilder} para telas adaptáveis.
    \end{itemize}
  \end{resumo}
\end{frame}

\begin{frame}{Glossário (1/2)}
  \begin{glossario}
    UI declarativa & a interface é descrita como função do estado; o framework aplica as mudanças \\
    Widget & descrição imutável e barata de uma parte da interface \\
    \cod{Element} & instância do widget na árvore; guarda o estado e o ciclo de vida \\
    \cod{RenderObject} & objeto que mede, posiciona e pinta \\
    \cod{BuildContext} & posição do widget na árvore; dá acesso a tema, tela e navegação \\
    \emph{Rebuild} & nova execução de \cod{build} nos elementos marcados como sujos \\
    Restrições & limites mínimos e máximos de largura e altura (\cod{BoxConstraints}) \\
    \cod{Expanded} & filho de \cod{Row}/\cod{Column} que ocupa o espaço que sobra \\
    Chave (\cod{Key}) & identidade estável entre reconstruções (\cod{ValueKey}, \cod{GlobalKey}) \\
  \end{glossario}
\end{frame}

\begin{frame}{Glossário (2/2)}
  \begin{glossario}
    \cod{InheritedWidget} & disponibiliza dados para toda uma subárvore (\cod{Theme.of}) \\
    \emph{Token} de design & decisão visual nomeada pelo papel (\cod{primary}, \cod{titleMedium}) \\
    \cod{ColorScheme.fromSeed} & gera a paleta Material 3 a partir de uma cor base \\
    Pixel lógico & unidade independente da densidade (\cod{devicePixelRatio}) \\
    Classe de tamanho & classe de largura da janela: compacta, média ou expandida \\
    \emph{Jank} & quadro perdido por exceder o orçamento de tempo \\
    Árvore de semântica & descrição da tela usada por leitores de tela como o TalkBack \\
  \end{glossario}
\end{frame}

\begin{frame}{Questões de revisão}
  \begin{questoes}
    \item Qual a diferença entre uma interface imperativa e uma declarativa? Por que a declarativa evita telas inconsistentes?
    \item Explique a regra ``restrições descem, tamanhos sobem, o pai posiciona'' com um exemplo numérico.
    \item Que papel têm as árvores de widgets, de elementos e de objetos de renderização?
    \item O que acontece, passo a passo, depois de uma chamada a \cod{setState}?
    \item Por que o método \cod{build} não deve iniciar uma requisição de rede?
    \item Por que uma \cod{ListView} dentro de uma \cod{Column} gera o erro de altura ilimitada? Como resolver?
    \item Quando é necessário passar uma \cod{Key} aos itens de uma lista?
    \item Por que usar \cod{Theme.of(context).colorScheme.primary} em vez de uma cor fixa?
  \end{questoes}
\end{frame}

\begin{frame}{Exercícios}
  \begin{exercicio}[Exercício 1 --- Layout]
    Reproduza um cartão de medidor: título, valor grande, unidade menor ao lado e uma \cod{LinearProgressIndicator} colorida conforme a faixa de temperatura.
  \end{exercicio}
  \begin{exercicio}[Exercício 2 --- Erros de layout]
    Coloque uma \cod{ListView} dentro de uma \cod{Column} e uma \cod{Row} com textos longos. Leia as mensagens de erro, use o Layout Explorer e corrija.
  \end{exercicio}
  \begin{exercicio}[Exercício 3 --- Responsivo]
    Faça a lista de equipamentos virar grade em telas largas e teste no navegador redimensionando a janela.
  \end{exercicio}
  \begin{projeto}[LabMonitor --- Etapa 5]
    Implemente a tela \textbf{Equipamentos} (protótipo do Módulo 04): \cod{Scaffold}, \cod{AppBar}, FAB e \cod{ListView.builder} de \cod{ItemEquipamento}, com o tema azul-petróleo nos modos claro e escuro.
  \end{projeto}
\end{frame}

\begin{frame}{Para aprofundar}
  \small
  \begin{itemize}
    \item Introdução aos widgets: \url{https://docs.flutter.dev/ui}
    \item Catálogo de widgets: \url{https://docs.flutter.dev/ui/widgets}
    \item Entendendo restrições: \url{https://docs.flutter.dev/ui/layout/constraints}
    \item Layouts: \url{https://docs.flutter.dev/ui/layout}
    \item Material 3 no Flutter: \url{https://docs.flutter.dev/ui/design/material}
    \item Design adaptativo: \url{https://docs.flutter.dev/ui/adaptive-responsive}
  \end{itemize}
\end{frame}

\begin{frame}{Respostas das questões de revisão}
  \begin{respostas}
    \item No imperativo, o código altera componentes existentes e precisa sincronizá-los a cada evento; no declarativo, descreve-se a tela para cada estado e o framework aplica as diferenças. Há um só caminho do estado à tela, então nada fica esquecido.
    \item Na tela de 360 pixels lógicos, a \cod{Column} passa ao cartão largura de 0 a 360 e altura livre; o cartão escolhe 360 $\times$ 96; a coluna o põe em $y = 0$ e o próximo em $y = 96$.
    \item Widgets são descrições baratas, recriadas a cada mudança; elementos são as instâncias duradouras que guardam o \cod{State}; objetos de renderização medem, posicionam e pintam, e são reaproveitados.
    \item O elemento é marcado como sujo; no quadro seguinte seu \cod{build} roda de novo; filhos de mesmo tipo e chave são reaproveitados e atualizados; depois vêm layout e pintura só do que mudou.
    \item Porque \cod{build} pode rodar muitas vezes por segundo: a requisição se repetiria a cada reconstrução. Ela pertence ao \cod{initState}, ao gerenciador de estado ou ao repositório (Módulos 07 e 08).
    \item A \cod{Column} oferece altura ilimitada ao filho e a lista tenta ocupar toda a altura disponível. Envolva a lista em \cod{Expanded} ou dê a ela uma altura definida.
    \item Quando itens com estado próprio podem ser inseridos, removidos ou reordenados: a \cod{ValueKey(id)} mantém o estado junto do item certo.
    \item O \emph{token} se adapta ao tema escuro, à cor dinâmica e a mudanças da identidade visual; a cor fixa, não.
  \end{respostas}
\end{frame}

\encerramento{Módulo 06 --- Dart assíncrono: Future, async/await, Stream e Isolates.}

\end{document}
